% Options for packages loaded elsewhere
\PassOptionsToPackage{unicode}{hyperref}
\PassOptionsToPackage{hyphens}{url}
\PassOptionsToPackage{table}{xcolor}
%
\documentclass[
  11pt,
]{report}
\usepackage{amsmath,amssymb}
\usepackage{iftex}
\ifPDFTeX
  \usepackage[T1]{fontenc}
  \usepackage[utf8]{inputenc}
  \usepackage{textcomp} % provide euro and other symbols
\else % if luatex or xetex
  % pdfLaTeX-only repository policy: the original unicode-math/fontspec
  % branch is unreachable under pdfLaTeX and is replaced with a guard.
  \PackageError{cpython-suite}{Compile with pdfLaTeX only}{}
\fi
\usepackage{lmodern}
\ifPDFTeX\else
  % xetex/luatex font selection
\fi
% Use upquote if available, for straight quotes in verbatim environments
\IfFileExists{upquote.sty}{\usepackage{upquote}}{}
\IfFileExists{microtype.sty}{% use microtype if available
  \usepackage[]{microtype}
  \UseMicrotypeSet[protrusion]{basicmath} % disable protrusion for tt fonts
}{}
\makeatletter
\@ifundefined{KOMAClassName}{% if non-KOMA class
  \IfFileExists{parskip.sty}{%
    \usepackage{parskip}
  }{% else
    \setlength{\parindent}{0pt}
    \setlength{\parskip}{6pt plus 2pt minus 1pt}}
}{% if KOMA class
  \KOMAoptions{parskip=half}}
\makeatother
\usepackage{xcolor}
\usepackage{listings}
\newcommand{\passthrough}[1]{#1}
\lstset{defaultdialect=[5.3]Lua}
\lstset{defaultdialect=[x86masm]Assembler}
\usepackage{longtable,booktabs,array}
\usepackage{calc} % for calculating minipage widths
% Correct order of tables after \paragraph or \subparagraph
\usepackage{etoolbox}
\makeatletter
\patchcmd\longtable{\par}{\if@noskipsec\mbox{}\fi\par}{}{}
\makeatother
% Allow footnotes in longtable head/foot
\IfFileExists{footnotehyper.sty}{\usepackage{footnotehyper}}{\usepackage{footnote}}
\makesavenoteenv{longtable}
\setlength{\emergencystretch}{3em} % prevent overfull lines
\providecommand{\tightlist}{%
  \setlength{\itemsep}{0pt}\setlength{\parskip}{0pt}}
\setcounter{secnumdepth}{5}
% CPython Engineering Suite shared presentation layer -- pdfLaTeX compatible
\usepackage[T1]{fontenc}
\usepackage{lmodern}
\usepackage{microtype}
\usepackage{geometry}
\geometry{letterpaper,margin=0.91in,headheight=16pt,headsep=19pt}
\usepackage{xcolor} % [table] passed at top of file (option-clash fix)
\definecolor{CPNavy}{HTML}{13283B}
\definecolor{CPBlue}{HTML}{176BA0}
\definecolor{CPGold}{HTML}{D8A73C}
\definecolor{CPMuted}{HTML}{566979}
\definecolor{CPBg}{HTML}{EFF5F8}
\usepackage{titlesec}
\titleformat{\chapter}[display]{\normalfont\sffamily\bfseries\color{CPNavy}}{\large\color{CPBlue}CHAPTER \thechapter}{9pt}{\Huge}[\vspace{6pt}\color{CPGold}\titlerule]
\titleformat{\section}{\Large\bfseries\sffamily\color{CPBlue}}{\thesection}{0.7em}{}
\titleformat{\subsection}{\large\bfseries\sffamily\color{CPNavy}}{\thesubsection}{0.7em}{}
\usepackage{fancyhdr}
\pagestyle{fancy}
\fancyhf{}
\fancyhead[L]{\footnotesize\color{CPMuted}\sffamily CPYTHON ENGINEERING SUITE}
\fancyhead[R]{\footnotesize\color{CPMuted}\sffamily 3.16.0a1+dev SNAPSHOT}
\fancyfoot[C]{\footnotesize\color{CPMuted}\thepage}
\renewcommand{\headrulewidth}{0.3pt}
\AtBeginDocument{\urlstyle{tt}}
\lstset{basicstyle=\ttfamily\small,columns=fullflexible,keepspaces=true,breaklines=true,breakatwhitespace=false,showstringspaces=false,frame=single,framesep=5pt,rulecolor=\color{CPBg},backgroundcolor=\color{CPBg},xleftmargin=3pt,xrightmargin=3pt}
\usepackage{tcolorbox}
\tcbuselibrary{breakable,skins}
\newtcolorbox{sourcebox}[1][]{enhanced,breakable,sharp corners,boxrule=0.4pt,colback=CPBg,colframe=CPBlue,borderline west={3pt}{0pt}{CPBlue},left=10pt,right=8pt,top=6pt,bottom=6pt,#1}
\newtcolorbox{cautionbox}[1][]{enhanced,breakable,sharp corners,boxrule=0.4pt,colback=CPBg,colframe=CPGold,borderline west={3pt}{0pt}{CPGold},left=10pt,right=8pt,top=6pt,bottom=6pt,#1}
\usepackage{tikz}
\usetikzlibrary{arrows.meta,positioning,shapes.geometric,calc}
\usepackage{enumitem}
\setlist{itemsep=2pt,topsep=4pt}
\usepackage{booktabs}
\usepackage{longtable}
\AtBeginEnvironment{longtable}{\small}
\usepackage{array}
\setlength{\parskip}{5pt}
\setlength{\parindent}{0pt}
\setcounter{tocdepth}{2}
\usepackage{hyperref}
\hypersetup{colorlinks=true,linkcolor=CPBlue,urlcolor=CPBlue,citecolor=CPBlue,pdfauthor={CPython Engineering Documentation Suite}}
\makeatletter
\renewcommand{\maketitle}{%
  \begin{titlepage}\pagecolor{CPNavy}\color{white}\sffamily
  \vspace*{0.8in}
  {\color{CPGold}\rule{2.1in}{4pt}}\\[16pt]
  {\Large\bfseries CPYTHON ENGINEERING DOCUMENTATION SUITE}\\[24pt]
  {\Huge\bfseries \@title\par}\vspace{19pt}
  \vfill
  {\normalsize Source baseline: CPython 3.16.0 alpha 1+dev repository snapshot}\par
  {\normalsize Prepared 16 October 2026}\par
  {\small\color{white!70} Source-traced engineering commentary; not an official Python specification.}\vspace{0.65in}
  \end{titlepage}\nopagecolor\color{black}
  }
\makeatother
\ifLuaTeX
  \usepackage{selnolig}  % disable illegal ligatures
\fi
\usepackage{bookmark}
\IfFileExists{xurl.sty}{\usepackage{xurl}}{} % add URL line breaks if available
\urlstyle{same}
\hypersetup{
  pdftitle={CPython Compiler and Execution Pipeline},
  hidelinks,
  pdfcreator={LaTeX via pandoc}}

\title{CPython Compiler and Execution Pipeline}
\usepackage{etoolbox}
\makeatletter
\providecommand{\subtitle}[1]{% add subtitle to \maketitle
  \apptocmd{\@title}{\par {\large #1 \par}}{}{}
}
\makeatother
\subtitle{Volume 04 \textbar{} Lexing, PEG parsing, AST, instruction
generation, CFG and bytecode}
\author{}
\date{}

\begin{document}
\maketitle

{
\setcounter{tocdepth}{1}
\tableofcontents
}
\chapter{Architecture and stage
contract}\label{architecture-and-stage-contract}

\begin{sourcebox}
\textbf{Evidence classification.} \textbf{SOURCE} = statements supported by material present in the uploaded snapshot; \textbf{IMPLEMENTATION} = CPython-specific internals, not guaranteed Python semantics; \textbf{ENGINEERING} = original recommendations, examples, checks, and exercises. The source baseline is \textbf{3.16.0 alpha 1+dev}; missing repository file bodies are noted as evidence limitations.
\end{sourcebox}

\path{InternalDocs/compiler.md} names five
source-to-bytecode stages: tokenization
(\path{Parser/lexer/} and
\path{Parser/tokenizer/}), parsing
(\path{Parser/parser.c}) into an abstract syntax
tree, instruction-sequence generation
(\path{Python/compile.c}), control-flow graph
construction/optimization (\path{Python/flowgraph.c})
and bytecode emission (\path{Python/assemble.c}). The
full bodies of some named generator and compiler sources were not
included in the uploaded text export, so this guide uses the internal
design documents to ground the stage architecture and avoids pretending
to have audited their missing code.

\textbf{IMPLEMENTATION boundary:} The compiler's AST representation, CFG
passes and emitted bytecode are CPython-specific. User programs should
rely on the Python language reference rather than assuming opcode or
code-object storage compatibility across releases.

\begin{longtable}[]{@{}
  >{\raggedright\arraybackslash}p{(\columnwidth - 6\tabcolsep) * \real{0.2500}}
  >{\raggedright\arraybackslash}p{(\columnwidth - 6\tabcolsep) * \real{0.2500}}
  >{\raggedright\arraybackslash}p{(\columnwidth - 6\tabcolsep) * \real{0.2500}}
  >{\raggedright\arraybackslash}p{(\columnwidth - 6\tabcolsep) * \real{0.2500}}@{}}
\toprule\noalign{}
\begin{minipage}[b]{\linewidth}\raggedright
Stage
\end{minipage} & \begin{minipage}[b]{\linewidth}\raggedright
Principal input
\end{minipage} & \begin{minipage}[b]{\linewidth}\raggedright
Principal output
\end{minipage} & \begin{minipage}[b]{\linewidth}\raggedright
Validation target
\end{minipage} \\
\midrule\noalign{}
\endhead
\bottomrule\noalign{}
\endlastfoot
Tokenization & Source text and encoding & Token stream & Correct
boundaries, indentation and literals \\
Parsing & Tokens and grammar & AST & Syntax accepted/rejected
correctly \\
AST processing & AST and symbol information & Instruction sequence &
Correct binding and semantic lowering \\
Optimization & Instructions and control flow & Optimized graph &
Semantics and exception paths retained \\
Assembly & Graph/labels/metadata & Code object bytecode & Jump offsets,
metadata, execution behavior \\
\end{longtable}

\chapter{Lexical analysis and parser
generation}\label{lexical-analysis-and-parser-generation}

\section{Tokenization}\label{tokenization}

Lexing separates physical source content into syntactic units and tracks
indentation/dedentation. Whitespace inside strings, comments, explicit
continuation and bracketed expressions behave differently. Consult
\path{Doc/reference/lexical_analysis.rst} for
user-facing rules and \path{InternalDocs/parser.md}
for how the generated parser consumes tokens.

\begin{lstlisting}[language=Python]
import io
import tokenize

source = b"total = (2 +\n         3)\n"
for token in tokenize.tokenize(io.BytesIO(source).readline):
    print(tokenize.tok_name[token.type], repr(token.string))
\end{lstlisting}

This program is a tokenization demonstration, not proof that public
tokenizer APIs exactly mirror all internal implementation choices.

\section{PEG semantics}\label{peg-semantics}

The parser is PEG-based;
\path{InternalDocs/parser.md} explains that
alternatives are \textbf{ordered}. A rule
\passthrough{\lstinline!A | B!} attempts \passthrough{\lstinline!A!}
first, and the grammar is not simply an unordered list of equivalent
context-free productions. Lookaheads, commitment/cut behavior and
invalid syntax recovery are detailed in
\path{InternalDocs/parser.md} and
\path{InternalDocs/changing_grammar.md}. The
formatted grammar described by
\path{Doc/reference/grammar.rst} omits details
related to code generation and error recovery.

\textbf{ENGINEERING parser-change checklist:} define syntax and expected
semantics; update grammar; check ambiguous-looking ordered alternatives;
ensure error messages remain useful; add positive and negative grammar
tests; update AST/compiler consumers; regenerate derived parser
artifacts with repository tooling; inspect formatting and docs.

\chapter{Abstract syntax trees and semantic
lowering}\label{abstract-syntax-trees-and-semantic-lowering}

\section{The AST is not source text}\label{the-ast-is-not-source-text}

\passthrough{\lstinline!ast.parse!} produces an AST that can be explored
and transformed. Source trivia such as comments and whitespace is not
generally preserved in the AST. An AST transformation must preserve
expected node location data, context
(\passthrough{\lstinline!Load!}/\passthrough{\lstinline!Store!}) and
semantics when compiled. Altering a tree without understanding scope,
evaluation order or exception behavior is unsafe.

\begin{lstlisting}[language=Python]
import ast

source = "result = amount * 2"
tree = ast.parse(source, mode="exec")
print(ast.dump(tree, indent=2))
code = compile(tree, filename="example.py", mode="exec")
namespace = {"amount": 4}
exec(code, namespace)
assert namespace["result"] == 8
\end{lstlisting}

\section{Symbol binding and scope}\label{symbol-binding-and-scope}

Symbol analysis distinguishes local, global, nonlocal, parameter, and
free-variable relationships. Python's execution model defines scope
semantics, while the compiler resolves them into operations. A closure
requires consistent handling of values captured from an enclosing
lexical environment; an optimizer cannot change an observable binding
relationship.

\begin{lstlisting}[language=Python]
import symtable

source = '''
def outer(a):
    scale = 3
    def inner(b):
        return a * scale + b
    return inner
'''
root = symtable.symtable(source, "example.py", "exec")
for child in root.get_children():
    print(child.get_name(), child.get_type())
\end{lstlisting}

The \passthrough{\lstinline!symtable!} module exposes compiler-related
analysis but should not be mistaken for a permanent representation of
every compiler internal.

\chapter{Instruction sequences, CFG and
optimization}\label{instruction-sequences-cfg-and-optimization}

\section{Flow graph invariants}\label{flow-graph-invariants}

Control flow represents normal edges, branches, exceptional exits, loop
back-edges and handler transitions. Optimization may remove unreachable
regions, simplify computations or rework jumps \textbf{only} when it
preserves observable behavior. Exception behavior, debugging line
information, cleanup and generator suspension create correctness
constraints beyond a simple arithmetic output check.

\textbf{ENGINEERING proof obligations for each transform:}

\begin{enumerate}
\def\labelenumi{\arabic{enumi}.}
\tightlist
\item
  The entry state is compatible with the transformed sequence.
\item
  Every reachable exit observes equivalent Python effects and
  return/exception results.
\item
  Required stack effects and ownership transitions remain balanced.
\item
  Exceptions and finally/cleanup routes retain their intended meaning.
\item
  Debug/trace information remains as accurate as required by the
  relevant interface.
\end{enumerate}

\section{Bytecode and assembly}\label{bytecode-and-assembly}

\path{InternalDocs/compiler.md} attributes final
code-object construction to the assembly stage.
\path{InternalDocs/code_objects.md} explains code
objects and associated execution metadata. Jump offsets and opcode
numbering are implementation details, not data formats for external
storage. \passthrough{\lstinline!dis!} is a practical observer for an
interpreter build and \passthrough{\lstinline!compile()!} permits
compilation in \passthrough{\lstinline!exec!},
\passthrough{\lstinline!eval!} or \passthrough{\lstinline!single!}
modes, each with different source expectations.

\begin{lstlisting}[language=Python]
import dis

source = '''
def absdiff(a, b):
    if a >= b:
        return a - b
    return b - a
'''
code = compile(source, "absdiff.py", "exec")
namespace = {}
exec(code, namespace)
dis.dis(namespace["absdiff"])
assert namespace["absdiff"](7, 2) == 5
\end{lstlisting}

\chapter{Exceptions, generators and special
lowering}\label{exceptions-generators-and-special-lowering}

Exception handlers depend on compilation metadata and runtime unwinding;
see \path{InternalDocs/exception_handling.md}.
Generators transform function activation into resumable execution state;
see \path{InternalDocs/generators.md}.
\path{InternalDocs/inlined_comprehensions.md}
describes compiler handling of inlined comprehensions in relevant cases.
These features expose why source-to-opcode translation is not a one-line
mapping.

\textbf{Invariant exercise:} Create a function containing a loop, a
comprehension, a \passthrough{\lstinline!try/finally!}, and a generator.
Inspect its code objects and annotate which syntactic constructs
introduce branches, nested code objects or suspension points in the
pinned interpreter.

\chapter{Compiler engineering
workflow}\label{compiler-engineering-workflow}

\section{Before changing grammar or
lowering}\label{before-changing-grammar-or-lowering}

Record the full version and branch, run a clean test baseline, and
identify generated files.
\path{InternalDocs/changing_grammar.md} and
\path{InternalDocs/parser.md} should guide grammar
regeneration; do not manually patch generated parser C code as if it
were authoritative source. When introducing a language feature, consider
syntax diagnostics, AST consumers, compiler lowering, docs,
debugger/profiler integration, deprecations, and downstream tooling
compatibility.

\section{A focused testing pyramid}\label{a-focused-testing-pyramid}

\begin{longtable}[]{@{}
  >{\raggedright\arraybackslash}p{(\columnwidth - 4\tabcolsep) * \real{0.3333}}
  >{\raggedright\arraybackslash}p{(\columnwidth - 4\tabcolsep) * \real{0.3333}}
  >{\raggedright\arraybackslash}p{(\columnwidth - 4\tabcolsep) * \real{0.3333}}@{}}
\toprule\noalign{}
\begin{minipage}[b]{\linewidth}\raggedright
Layer
\end{minipage} & \begin{minipage}[b]{\linewidth}\raggedright
Suggested test
\end{minipage} & \begin{minipage}[b]{\linewidth}\raggedright
Primary failure caught
\end{minipage} \\
\midrule\noalign{}
\endhead
\bottomrule\noalign{}
\endlastfoot
Tokenization & Token boundaries and indentation cases & Lexical
regression \\
Grammar & Accepted and rejected program snippets & Syntax ambiguity or
incorrect error \\
AST & Shape, location metadata, context flags & Semantic mapping
regression \\
Compiler & Execute code under multiple scopes & Binding/control-flow
errors \\
Runtime & Exceptions, closures, async/generator cases & Unwinding or
suspension defects \\
Differential & Compare stable and changed builds where semantics should
match & Behavioral regression \\
Performance & Repeated fixed workload with variance & Compilation or
execution cost changes \\
\end{longtable}

\chapter{Laboratories}\label{laboratories}

\section{Lab A: Reconstruct a source-to-bytecode
trace}\label{lab-a-reconstruct-a-source-to-bytecode-trace}

Define a nontrivial \passthrough{\lstinline!for!} loop that
conditionally accumulates values. Save the source,
\passthrough{\lstinline!tokenize!} output,
\passthrough{\lstinline!ast.dump!} result,
\passthrough{\lstinline!symtable!} output and
\passthrough{\lstinline!dis!} output. Write one paragraph per stage
explaining what information is created or discarded. Include Python
build identification.

\section{Lab B: AST rewriting with
invariants}\label{lab-b-ast-rewriting-with-invariants}

Implement an \path{ast.NodeTransformer} to change
only integer literal \passthrough{\lstinline!1!} into
\passthrough{\lstinline!2!} in a controlled expression. Call
\path{ast.fix_missing_locations}, compile the
transformed AST and test with \passthrough{\lstinline!unittest!}. Then
explain why such a transformation might be incorrect for arbitrary
programs.

\begin{lstlisting}[language=Python]
import ast

class ReplaceOne(ast.NodeTransformer):
    def visit_Constant(self, node):
        if type(node.value) is int and node.value == 1:
            return ast.copy_location(ast.Constant(value=2), node)
        return node

root = ast.parse("result = 1 + 3")
root = ast.fix_missing_locations(ReplaceOne().visit(root))
state = {}
exec(compile(root, "<transformed>", "exec"), state)
assert state["result"] == 5
\end{lstlisting}

\section{Lab C: Scope test matrix}\label{lab-c-scope-test-matrix}

Create local, closure, global and \passthrough{\lstinline!nonlocal!}
examples. Predict resolution before running them. If any outcome
differs, correct the model using
\path{Doc/reference/executionmodel.rst}.

\section{Lab D: CFG and exception
reasoning}\label{lab-d-cfg-and-exception-reasoning}

Compare \passthrough{\lstinline!return!} from inside
\passthrough{\lstinline!try!} with an exception raised in
\passthrough{\lstinline!finally!}. Describe observable precedence and
identify why an optimizer cannot remove the cleanup path.

\textbf{Acceptance:} each lab includes executable code, known expected
output, source paths, and a discussion that distinguishes the public
language contract from implementation artifacts.

\chapter{Source concordance and known
limits}\label{source-concordance-and-known-limits}

\begin{itemize}
\tightlist
\item
  \textbf{SOURCE:} \path{InternalDocs/compiler.md}
  (five-stage compiler design), \passthrough{\lstinline!parser.md!},
  \path{changing_grammar.md},
  \path{code_objects.md},
  \path{exception_handling.md},
  \path{generators.md},
  \path{inlined_comprehensions.md}.
\item
  \textbf{SOURCE:}
  \path{Doc/reference/lexical_analysis.rst},
  \passthrough{\lstinline!grammar.rst!},
  \path{executionmodel.rst}; module documentation for
  \passthrough{\lstinline!ast!}, \passthrough{\lstinline!dis!},
  \passthrough{\lstinline!symtable!} where present.
\item
  \textbf{BODY PRESENT:} \path{Python/assemble.c},
  \path{Python/instruction_sequence.c},
  \passthrough{\lstinline!Python/ast.c!},
  \path{Python/ast_unparse.c}.
\item
  \textbf{REFERENCE / TREE ONLY in this supplied export:} some large
  compiler/parser files (\path{Grammar/python.gram},
  \path{Parser/parser.c},
  \path{Python/compile.c},
  \path{Python/flowgraph.c}). Claims about their
  detailed code bodies must be independently checked in a complete
  checkout.
\end{itemize}

The examples and test matrices are original teaching material. They are
not quotations from the Python specification.

\end{document}
